\Askhsh{22130_SOLUTION}{1}
\begin{ylh}{1}
    \item[ΛΥΣΗ]
\end{ylh}
\begin{wrapfigure}{r}{5.5cm}
    \centering
    \begin{tikzpicture}[scale=0.9]
        % Define vertices of the right triangle
        \tkzDefPoint(0,0){A}
        \tkzDefPoint(0,3.5){B}
        \tkzDefPoint(4,0){G}

        % Find the circumcenter (midpoint of hypotenuse)
        \tkzDefMidPoint(B,G)\tkzGetPoint{D}

        % Draw the circle
        \tkzDrawCircle[color=black](D,A)

        % Draw the triangle and radii
        \tkzDrawPolygon[color=black](A,B,G)
        \tkzDrawSegments[color=black](D,A D,B D,G)

        % Mark the right angle
        \tkzMarkRightAngle[color=black](B,A,G)

        % Label points
        \tkzLabelPoint[below left](A){A}
        \tkzLabelPoint[above left](B){B}
        \tkzLabelPoint[below right](G){$\Gamma$}
        \tkzLabelPoint[above right](D){$\Delta$}

        % Label sides and radii
        \tkzLabelSegment[left, pos=0.5](A,B){$\gamma$}
        \tkzLabelSegment[below, pos=0.5](A,G){$\beta$}
        \tkzLabelSegment[above, sloped, pos=0.6](D,B){R}
        \tkzLabelSegment[above, sloped, pos=0.6](D,G){R}
        \tkzLabelSegment[below right, sloped, pos=0.5](D,A){R}

        % Label angle alpha. Based on the diagram, this appears to be angle AΔΓ
        \tkzMarkAngle[size=0.6, color=black](G,D,A)
        \tkzLabelAngle[pos=0.8, color=black](G,D,A){$\alpha$}
    \end{tikzpicture}
\end{wrapfigure}
\begin{alist}
    \item Στο σχήμα, το ορθογώνιο τρίγωνο ΑΒΓ είναι εγγεγραμμένο σε κύκλο ακτίνας R. Αφού η εγγεγραμμένη γωνία $\hat{A}$ είναι ορθή, τότε θα βαίνει σε ημικύκλιο. Επομένως, η υποτείνουσα ΒΓ του τριγώνου είναι διάμετρος του κύκλου. Άρα ΒΓ = α = 2R.
    \item Από το Πυθαγόρειο θεώρημα στο ορθογώνιο τρίγωνο ΑΒΓ έχουμε:
    \[ ΒΓ^2 = ΑΓ^2 + ΑΒ^2 \text{ ή } α^2 = β^2 + γ^2 \]
\end{alist}
Οπότε:
\[ α^2 + β^2 + γ^2 = α^2 + α^2 = 2α^2 = 2 \cdot (2R)^2 = 2 \cdot 4R^2 = 8R^2 \]